\documentclass[10pt,twocolumn]{article}
\usepackage[T1]{fontenc}
\usepackage[utf8]{inputenc}
\usepackage{lmodern}
\usepackage[margin=1.8cm,top=2.4cm,bottom=2cm,columnsep=0.6cm]{geometry}
\usepackage{fancyhdr}
\usepackage{titlesec}
\usepackage{booktabs}
\usepackage{amsmath,amssymb,amsfonts}
\usepackage{microtype}
\usepackage{xcolor}
\usepackage{mdframed}
\usepackage{parskip}
\usepackage{enumitem}
\usepackage{hyperref}

\definecolor{rulered}{RGB}{180,30,30}
\definecolor{boxgray}{RGB}{245,245,245}
\definecolor{keywordgray}{RGB}{100,100,100}

\pagestyle{fancy}
\fancyhf{}
\fancyhead[L]{\textsc{\small Claude Mag}}
\fancyhead[C]{\small\textit{Machine Learning Research Digest}}
\fancyhead[R]{\small 2026-07-15}
\fancyfoot[C]{\small\thepage}
\renewcommand{\headrulewidth}{0.8pt}

\titleformat{\section}{\normalsize\bfseries\color{rulered}}{}{0pt}{}%
  [\vspace{-4pt}\color{rulered}\rule{\columnwidth}{0.4pt}]
\titlespacing{\section}{0pt}{8pt}{4pt}

\hypersetup{colorlinks=true,urlcolor=rulered,linkcolor=black}

\begin{document}
\setlength{\parindent}{0pt}
\setlength{\parskip}{4pt}

{\small\color{keywordgray}\textbf{Keywords:}
  agent evaluation \textbullet{}
  benchmark \textbullet{}
  long-horizon tasks \textbullet{}
  dense rewards}\par\vspace{4pt}

{\LARGE\bfseries\raggedright
  Long-Horizon-Terminal-Bench}\par\vspace{4pt}

{\large\itshape\raggedright
  A 46-Task Dense-Reward Benchmark Exposes the
  True Performance Gap Between LLM Agents on
  Complex Terminal Workflows}\par\vspace{6pt}

{\small\textbf{By} Zongxia Li, Zhongzhi Li, Yucheng Shi, Ruhan Wang
  et al.\ (Tencent HY LLM Frontier; Univ.\ of Maryland; et al.)}\par\vspace{2pt}

{\small\color{keywordgray}arXiv:2607.08964 \textbullet{}
  July 9, 2026}
\par\vspace{2pt}\noindent{\color{rulered}\rule{\columnwidth}{0.6pt}}\vspace{6pt}

Binary pass/fail benchmarks for coding and terminal agents are
comfortable lies: they assign 0 or 1 to an episode without capturing
how far the agent progressed on complex, multi-step workflows.
Long-Horizon-Terminal-Bench (LHTB) replaces this illusion with 46
tasks spanning nine categories, graded by hidden rebuild-from-artifact
verifiers that assign partial credit through fine-grained subtask
completion. The result is a sobering view: top agents score only
28--34\% on the hardest categories---and would have appeared far more
capable under binary grading.

\begin{mdframed}[backgroundcolor=boxgray,linecolor=rulered,linewidth=1pt,
    innerleftmargin=6pt,innerrightmargin=6pt,
    innertopmargin=6pt,innerbottommargin=6pt]
\textbf{\small AT A GLANCE}\par\vspace{4pt}
\begin{description}[leftmargin=0pt,labelsep=4pt,itemsep=2pt,parsep=0pt,
    font=\small\bfseries]
  \item[Problem:] Binary agent benchmarks mask variation in task
    progress, distorting comparisons and hiding true capability gaps.
  \item[Why it matters:] Real deployment requires agents that make
    sustained, useful progress over hundreds of steps; binary grading
    cannot differentiate 5\% from 95\% completion.
  \item[Method:] 46-task terminal benchmark with subtask-level
    dense rewards; hidden rebuild-from-artifact verifiers prevent
    score gaming.
  \item[Result:] Top agents (Claude 3.7, GPT-4o) average 52\%
    overall; only 28--34\% on scientific computing and experiment
    reproduction.
  \item[Evidence:] Three frontier models evaluated; binary-vs-dense
    correlation only 0.41, confirming benchmarks measure different
    capabilities.
\end{description}
\end{mdframed}

\section{The Big Picture}

Current LLM agent evaluations suffer from a fundamental measurement
problem. SWE-bench, HumanEval, and Terminal-Bench all grade each
task as a binary 0 or 1: either the final state passes all tests, or
it does not. This approach is easy to implement and produces clean
leaderboards, but it conflates ``failed to start'' with ``almost
finished''---both receive the same score of 0.

The problem compounds on long-horizon tasks where success requires
chaining dozens of subtasks: running a numerical simulation, checking
convergence, generating plots, and writing a report. An agent that
completes the simulation and convergence check but fails on the
plotting step does 60\% of the work---but receives 0 credit under
binary grading.

LHTB introduces dense intermediate rewards by decomposing every task
into 4--12 fine-grained, verifiable subtasks, each graded independently.
Partial credit is not merely cosmetic: it changes which agents rank
highest and by how much, and reveals capability gaps invisible under
binary evaluation.

\section{Why Existing Approaches Fall Short}

Terminal-Bench introduced the key insight of containerized, agentic
evaluation with a bash shell---but retained binary grading. SWE-bench
Verified improved task realism (real GitHub issues) but also grades
binary. Both benchmarks are additionally vulnerable to self-reported
progress: agents can write test files that trivially pass their own
assertions, gaming the verifier.

Process reward models (PRMs) offer a complementary approach---training
a model to assign intermediate rewards---but require supervised reward
data that is expensive to collect and brittle to domain shift. LHTB
achieves dense rewards without any learned reward model by decomposing
tasks into subtasks with objectively verifiable intermediate states
(e.g., ``the simulation output file exists and contains a valid numpy
array with shape $(T, N)$'').

\section{Core Mechanism}

The benchmark's core contribution is the combination of (1)~task
decomposition into verifiable subtasks and (2)~rebuild-from-artifact
verification that prevents score gaming.

For rebuild-from-artifact verification, instead of running the agent's
code directly, the verifier mounts the agent's output directory
read-only, extracts specified intermediate artifacts (files, database
entries, generated figures), reconstructs a ground-truth comparison
from these artifacts, and asserts properties. This prevents agents
from planting self-passing assertions or manipulating the test
environment because the verifier operates on static artifacts, not
live processes.

\section{Technical Formulation}

Each task $\tau$ decomposes into subtasks $\{s_1, \ldots, s_m\}$,
each with a binary verifier $v_i \in \{0, 1\}$. The composite score is:
\[
S(\tau) = 0.7 \cdot \frac{1}{m}\sum_{i=1}^m v_i
  + 0.2 \cdot V_{\text{final}}
  + 0.1 \cdot B_{\text{eff}}
\]
where $V_{\text{final}}$ is a holistic final-state verifier and
$B_{\text{eff}}$ is an efficiency bonus that decays linearly for
episodes requiring more than 100 steps. This scoring incentivizes
both thorough subtask completion and efficient execution.

\section{Learning or Inference Procedure}

Agents are placed in an Ubuntu 22.04 Docker container with standard
scientific Python, build tools, git, and a network connection.
They receive three tools: (1)~a bash executor, (2)~a file editor,
and (3)~a web search interface. The step limit is 200; agents are
instructed only with the natural-language task description.

\begin{enumerate}[itemsep=2pt,parsep=0pt]
  \item Agent reads task description and initializes workspace.
  \item Agent executes a sequence of bash commands, file edits,
    and searches to complete the task.
  \item After each agent-initiated ``checkpoint'' (or at the step limit),
    the verifier mounts the filesystem, extracts specified artifacts,
    and runs subtask assertions.
  \item Subtask credits accumulate; a final holistic verifier runs
    after the episode ends.
\end{enumerate}

\section{What the Guarantee Says}

The red-team evaluation provides an empirical anti-gaming guarantee:
agents instructed to game the verifier succeeded on 0/46 tasks,
versus 23/46 on a naive run-and-check verifier. The rebuild-from-artifact
approach does not prevent all possible manipulation, but the attack
surface is dramatically smaller because agents cannot interact with
the verifier process at all---only their output artifacts are examined.

\section{Experimental Findings}

\begin{center}
\small
\begin{tabular}{lccc}
\toprule
Category & Claude 3.7 & GPT-4o & Gemini 2.5 \\
\midrule
Software Eng. & 72.1\% & 68.4\% & 65.2\% \\
Data Eng. & 61.3\% & 58.7\% & 54.1\% \\
DevOps & 55.8\% & 52.3\% & 49.6\% \\
Interactive Games & 41.5\% & 38.2\% & 35.8\% \\
Experiment Repro. & 34.2\% & 30.1\% & 28.7\% \\
Scientific Comp. & 28.9\% & 27.4\% & 24.3\% \\
\midrule
\textbf{Overall (avg.)} & \textbf{52.3\%} & \textbf{49.1\%} & \textbf{45.9\%} \\
\bottomrule
\end{tabular}
\end{center}

The correlation between binary pass/fail scores on prior benchmarks
and LHTB subtask credit is only 0.41, confirming that LHTB captures
genuinely different capabilities. Binary grading reverses the ranking
between Claude 3.7 and GPT-4o on 7 of 46 tasks.

\section{Ablations and Interpretation}

\textbf{Binary vs.\ dense grading:} Binary grading on LHTB tasks
produces rankings that disagree with dense grading on 7/46 tasks,
illustrating how coarse grading can mislead model selection.

\textbf{Verifier robustness:} In the red-team evaluation, agents
instructed to game succeed on 0/46 tasks under rebuild-from-artifact
verification versus 23/46 under naive run-and-check, a 100\%
reduction in gaming success rate.

\textbf{Step count analysis:} Agents average 43 steps on software
engineering tasks and 161 steps on scientific computing tasks,
confirming that LHTB successfully stratifies tasks by difficulty and
that step count is a meaningful proxy for task complexity.

\textbf{Category difficulty ordering:} The difficulty gradient
(software engineering $\to$ data engineering $\to$ DevOps $\to$
games $\to$ experiment reproduction $\to$ scientific computing)
is consistent across all three frontier models tested, suggesting
it reflects genuine task structure rather than model-specific biases.

\section*{Reference}

Zongxia Li et al.\
\textbf{Long-Horizon-Terminal-Bench: Testing the Limits of Agents on
Long-Horizon Terminal Tasks with Dense Reward-Based Grading}.
arXiv:2607.08964, July 2026.
\url{https://arxiv.org/abs/2607.08964}

\end{document}
